\documentclass[10pt]{article}
\usepackage[T1]{fontenc}
\usepackage[utf8]{inputenc}
\usepackage[french]{babel}
\usepackage{amsmath,amssymb,mathrsfs,mathpazo}
\usepackage{geometry}
\geometry{a4paper,margin=2cm}
\usepackage{graphicx,array,tabularx,tabularray,colortbl,multirow,multicol,booktabs,enumitem,pifont,ulem,fancybox,diagbox,cancel}
\usepackage{pstricks,pstricks-add,pst-plot,pst-tree,pst-node,pst-text,pst-func,pst-math,pst-bspline,pst-3dplot,pst-eucl,pst-solides3d}
\let\pscalculate\undefined
\usepackage{pst-bezier}
\usepackage{tikz,pgfplots,tkz-tab}
\usetikzlibrary{arrows,arrows.meta,patterns,calc,positioning,trees,decorations.markings,shapes.geometric,shapes.misc,shapes.arrows,intersections,decorations.pathmorphing,decorations.pathreplacing,backgrounds,shadows,babel}
\pgfplotsset{compat=1.18}
\usepgfplotslibrary{fillbetween}
\usepackage[np]{numprint}
\usepackage{listings,makecell,pas-tableur,twoopt}
\let\barre\undefined
\usepackage[pstricks]{tablvar}
\definecolor{darkgreen}{RGB}{0,100,0}
\providecommand{\red}{\color{red}}
\providecommand{\blue}{\color{blue}}
\providecommand{\green}{\color{green}}
\providecommand{\black}{\color{black}}
\providecommand{\white}{\color{white}}
\providecommand{\gray}{\color{gray}}
\providecommand{\cyan}{\color{cyan}}
\providecommand{\magenta}{\color{magenta}}
\providecommand{\yellow}{\color{yellow}}
\providecommand{\violet}{\color{violet}}
\providecommand{\orange}{\color{orange}}
\providecommand{\euro}{\texteuro}
\providecommand{\EUR}{\texteuro}
\providecommand{\ieme}{\textsuperscript{e}}
\providecommand{\degres}{\ensuremath{^\circ}}
\providecommand{\degree}{\ensuremath{^\circ}}
\providecommand{\vv}[1]{\overrightarrow{#1}}
\providecommand{\Coord}[1]{#1}
\providecommand{\decosix}{}
\providecommand{\starredbullet}{\ensuremath{\bullet}}
\pagestyle{empty}
\setlength{\parindent}{0pt}
\usepackage{adjustbox}\DeclareUnicodeCharacter{2236}{:}\DeclareUnicodeCharacter{2190}{\ensuremath{\leftarrow}}
\providecommand{\R}{\mathbb{R}}
\providecommand{\N}{\mathbb{N}}
\providecommand{\Z}{\mathbb{Z}}
\providecommand{\Q}{\mathbb{Q}}
\providecommand{\D}{\mathbb{D}}
\providecommand{\C}{\mathbb{C}}
\providecommand{\e}{{\mathrm{e}}}
\providecommand{\vect}[1]{\overrightarrow{#1}}
\providecommand{\vectt}[1]{\overrightarrow{\mathrm{#1}}}
\providecommand{\vvt}[1]{\overrightarrow{\mathrm{#1}}}
\providecommand{\np}[2][]{{#2~#1}}
\providecommand{\vv}[1]{{\overrightarrow{#1}}}
\renewcommand{\d}{{\mathrm{d}}}
\providecommand{\degres}{{^\circ}}
\providecommand{\degre}{{^\circ}}
\providecommand{\blue}{\color{blue}}
\providecommand{\red}{\color{red}}
\providecommand{\enskip}{\quad}
\providecommand{\llbracket}{\lbrack\!\lbrack}
\providecommand{\rrbracket}{\rbrack\!\rbrack}
\providecommand{\euro}{{\text{€}}}
\providecommand{\no}{{\text{n°}}}
\providecommand{\interval}[2]{{[#1;#2]}}
\providecommand{\intervFO}[2]{{[#1;#2[}}
\providecommand{\dotfill}{\ldots}

\providecommand{\ds}{\displaystyle}
\providecommand{\textding}[1]{#1}
\providecommand{\Oij}{\(\mathrm{O};\vec{\imath},\vec{\jmath}\)}
\providecommand{\Oijk}{\(\mathrm{O};\vec{\imath},\vec{\jmath},\vec{k}\)}
\providecommand{\Ouv}{\(\mathrm{O};\vec{u},\vec{v}\)}
\providecommand{\intervFF}[2]{\ensuremath{[#1;#2]}}
\renewcommand{\euro}{\ifmmode\text{\texteuro}\else\texteuro\fi}
\DeclareUnicodeCharacter{00B0}{\ifmmode^{\circ}\else\textdegree\fi}
\DeclareUnicodeCharacter{20AC}{\ifmmode\text{\texteuro}\else\texteuro\fi}
% Énoncés : Éducation nationale. Transcription et corrigés : APMEP, Denis Vergès et contributeurs. Découpage et classement : CoopMaths, Sébastien Lozano, Éric Elter et contributeurs. Adaptation de la présentation : Maths Explorer.
% CC BY-SA : https://coopmaths.fr/www/about/
% https://coopmaths.fr/alea/static/bac/2021/tex/bac_2021_06_etrangers_9_cor.tex
\begin{document}
\noindent\textit{APMEP / CoopMaths. CC BY-SA. Présentation adaptée par Maths Explorer.}\par\medskip
\bigskip

\textbf{Partie A :}

\medskip















\begin{enumerate}
\item 
\starredbullet~Si le test est négatif on aura fait un test : $X_n = 1$ ;

\starredbullet~~Si le test est positif il faudra faire le test des $n$ personnes plus le test global : $X_n = 1 + n$.
\item 
$P\left(X_n = 1\right)$ est la probabilité que l'on ne fasse qu'un test : le test des $n$ personnes et que celui-ci soit négatif donc que les $n$ personnes ne soient pas malades.

La probabilité qu'une personne soit malade est égale à 0,05, donc qu'une personne soit saine $1 - 0,05 = 0,95$ et donc que $n$ personnes soient saines est égale à $0,95^n$. Donc $P\left(X_n = 1\right) = 0,95^n$.



On a donc $P\left(X_n = n + 1\right) = 1 - 0,95^n$
\begin{center}
\begin{tabularx}{0.6\linewidth}{|*{3}{>{\centering \arraybackslash}X|}}\hline
$x_i$						& 1 &$n + 1$\\ \hline
$P\left(X_n = x_i\right)$	&$0,95^n$	&$1 - 0,95^n$\\ \hline
\end{tabularx}
\end{center}

\item 

On a $E\left(X_n\right) = 1 \times 0,95^n + (n + 1) \times \left(1 - 0,95^n \right) = 0,95^n + (n + 1) -0,95^n(n + 1) = 0,95^n + n + 1 - n \times 0,95^n  - 0,95^n = n + 1 - n \times  0,95^n$.


Cette espérance représente le nombre moyen d'analyses à effectuer pour un échantillon $n$ personnes : cette espérance est voisine de $n$.
\end{enumerate}

\bigskip

\textbf{Partie B :}

\medskip

\begin{enumerate}
\item 

$f(x) = \ln (x) + x \ln (0,95)$.

$f$ est une somme de fonctions dérivables sur $[20~;~ +\infty[$ et sur cet intervalle :

$f'(x) = \dfrac{1}{x} + \ln (0,95)$.



Or $20 \leqslant x \Rightarrow \dfrac{1}{x} \leqslant \dfrac{1}{20} \Rightarrow \dfrac{1}{x} + \ln 0,95\leqslant \dfrac{1}{20} + \ln 0,95$.

Or $\dfrac{1}{20} + \ln 0,95 \approx - 0,001$ ; il en résulte que $f'(x) \leqslant 0$ : la fonction $f$ est décroissante sur $[20~;~ +\infty[$.
\item On rappelle que $\displaystyle\lim_{x \to + \infty} \dfrac{\ln x}{x} = 0$. 


On a puisque $x \ne 0$, \, $f(x) = x\left(\dfrac{\ln x}{x} + \ln 0,95 \right)$.

Puisque  $\displaystyle\lim_{x \to + \infty} \dfrac{\ln x}{x} = 0$, alors  $\displaystyle\lim_{x \to + \infty} \dfrac{\ln x}{x} + \ln 0,95 = \ln 0,95 < 0$.

Finalement par produit de limites puisque $\displaystyle\lim_{x \to + \infty} x = + \infty$ et que $\displaystyle\lim_{x \to + \infty} \dfrac{\ln x}{x} + \ln 0,95 = 0,95 < 0$, on a 

$\displaystyle\lim_{x \to + \infty}f(x) = - \infty$.
\item 
La question précédente montre que $f$ est décroissante de $f(20) = \ln 20 + 20\ln 0,95 \approx 1,97$ à moins l'infini.

$f$ est continue car dérivable sur $[20~;~ +\infty[$ et prend ses valeurs dans les l'intervalle 
$]- \infty~;~f(20)]$\, avec $f(20) \approx 1,97$.

Comme $0 \in ]- \infty~;~f(20)]$, et que $f$ est décroissante sur cet intervalle  il existe donc un réel unique $\alpha$ tel que $f(\alpha) = 0$.

La calculatrice donne $f(87,0) \approx 0,003$ et $f(87,1) \approx \np{- 0,0005}$, donc 
$87,0 < \alpha 87,1$.

\item 
D'après la question précédente :

$f(x) > 0$ sur $[20~;~\alpha[$ et $f(x) < 0$ sur $]\alpha~;~+ \infty[$. 
\end{enumerate}

\bigskip

\textbf{Partie C :}

\medskip








On a $E\left(X_n\right)  < n \iff n + 1 - n \times 0,95^n < n \iff 1 < n \times 0,95^n$ ou en prenant le logarithme népérien $0 < \ln n + \ln \left(0,95^n\right) \iff 0 < \ln n + n \ln 0,95 \iff 0 < f(n) \iff f(n) > 0$.

Or on a vu à la fin de la partie B que la fonction $f$ est positive sur l'intervalle [20~;~87].

Conclusion : tester toutes les personnes conduira à moins d'analyses qu'avec la méthode 1 avec des échantillons de 20 à 87 personnes au maximum. Au delà il vaut mieux utiliser la première méthode.
\bigskip

\end{document}
